\documentclass[11pt]{article}
\usepackage{fontspec}
\usepackage{amsmath,amssymb,amsthm}
\usepackage[margin=2.6cm]{geometry}
\usepackage{booktabs}
\usepackage{hyperref}

\newtheorem{theorem}{Theorem}
\newtheorem{corollary}[theorem]{Corollary}

\title{Why the $27\times27$ cubic system has an analytic inverse:\\
       the commutant of $O_h$ on $\mathbb{R}^3\otimes\mathbb{R}^3\otimes\mathbb{R}^3$}
\author{}
\date{}

\begin{document}
\maketitle

\begin{abstract}
The dynamic single-site system for a cubic voxel is a $27\times27$ matrix $M$
carrying $19$ independent entries. Inverting it by brute force produces an
expression of no use to anybody. This note explains why that is the wrong thing
to do: $M$ commutes with the octahedral group acting on
$\mathbb{R}^3\otimes\mathbb{R}^3\otimes\mathbb{R}^3$, and Schur's lemma then
forces $M$ into block-diagonal form in a basis fixed by the \emph{symmetry
alone}. The parameters $a,\dots,t$ move the eigenvalues; they cannot move the
eigenvectors. Inversion therefore reduces to inverting a handful of small
blocks, analytically and in general moduli.
\end{abstract}

%======================================================================
\section{The space and the group}
%======================================================================

The unknowns are \emph{raw} derivatives --- $\partial_p\partial_q u_i$ with all
three indices free and no symmetrisation --- so the solution space is
\[
  V \;=\; \mathbb{R}^3\otimes\mathbb{R}^3\otimes\mathbb{R}^3,
  \qquad \dim V = 27 .
\]
This is worth stating plainly because it is \emph{not} the space used by the
Voigt-based $T_{27}$, which is $3\oplus6\oplus18$ (displacement, symmetric
strain, quadratic) and projects the antisymmetric part out. The two spaces have
the same dimension and different content.

A cube is invariant under the octahedral group $O_h$, realised concretely as the
$48$ signed permutation matrices --- every $3\times3$ matrix with exactly one
$\pm1$ in each row and column. Because $V$ is a threefold tensor power, each
$R\in O_h$ acts on it as
\[
  \rho(R) \;=\; R\otimes R\otimes R .
\]

\paragraph{Verified.} $M$ commutes with $\rho$ on the generators of $O_h$ --- a
$90^\circ$ rotation about $z$, a $120^\circ$ rotation about $(111)$, and the
mirror $z\mapsto-z$:
\[
  \rho(R)\,M\,\rho(R)^{-1} - M \;=\; 0
  \qquad\text{identically in } a,\dots,t .
\]
Checked symbolically, not numerically: the difference is the zero matrix as a
polynomial identity in the nineteen parameters, so it holds for every medium and
every frequency, not at a sampled point.

%======================================================================
\section{The consequence}
%======================================================================

\begin{theorem}[Schur]
Let $V=\bigoplus_\lambda \Gamma_\lambda\otimes\mathbb{C}^{m_\lambda}$ be the
isotypic decomposition of $V$ under $G$, where $\Gamma_\lambda$ is an irreducible
representation of dimension $d_\lambda$ occurring with multiplicity $m_\lambda$.
Then every $G$-equivariant operator $M$ has the form
\[
  M \;=\; \bigoplus_\lambda \; \mathrm{id}_{\Gamma_\lambda} \otimes\, B_\lambda ,
  \qquad B_\lambda \in \mathbb{C}^{m_\lambda\times m_\lambda},
\]
and the commutant algebra is $\mathrm{End}_G(V)\cong\bigoplus_\lambda
\mathbb{C}^{m_\lambda\times m_\lambda}$, of dimension $\sum_\lambda m_\lambda^2$.
\end{theorem}

The content for us is the placement of the tensor factors. $M$ acts as the
\emph{identity} on the internal index of each irreducible representation and as
an $m_\lambda\times m_\lambda$ matrix on the multiplicity space. Two things
follow, and they are what make the problem tractable.

\begin{corollary}[The basis is parameter-free]
A symmetry-adapted basis block-diagonalises $M$ for \emph{all} values of
$a,\dots,t$ simultaneously. The nineteen parameters enter only through the
entries of the small matrices $B_\lambda$; they cannot rotate the blocks,
because the blocks are determined by $O_h$ before $M$ is known.
\end{corollary}

\begin{corollary}[Inversion is small]
\[
  M^{-1} \;=\; \bigoplus_\lambda \;\mathrm{id}_{\Gamma_\lambda}\otimes\, B_\lambda^{-1},
\]
so inverting a $27\times27$ matrix costs exactly the inversion of the
$B_\lambda$, whose sizes are the \emph{multiplicities} $m_\lambda$, each result
reused $d_\lambda$ times. A brute-force symbolic inverse of $M$ recomputes the
same small quantities $27$ times over, in a basis that hides them.
\end{corollary}

%======================================================================
\section{What is measured}
%======================================================================

\begin{table}[h]
\centering
\begin{tabular}{lr}
\toprule
quantity & value \\
\midrule
$|O_h|$ & $48$ \\
$\dim\mathrm{End}_G(V)=\sum_\lambda m_\lambda^2
  =\frac{1}{48}\sum_{R}(\operatorname{tr}R)^6$ & $31$ \\
free parameters in $M$ & $19$ \\
generic eigenvalue degeneracies of $M$ & $\{1,\,2,\,3^{\times8}\}$ \\
\bottomrule
\end{tabular}
\end{table}

The commutant dimension is computed exactly from the character
$\chi(R)=(\operatorname{tr}R)^3$ of $\rho$, with no character table required,
since $\sum_\lambda m_\lambda^2 = \langle\chi,\chi\rangle$.

Two readings of this table matter.

\paragraph{$M$ does not fill its commutant.} Nineteen parameters inside a
$31$-dimensional algebra means $M$ carries $12$ further relations beyond
$O_h$-equivariance. Those are physics, not symmetry --- reciprocity and the
structure of the Green operator --- and they are worth identifying, because each
one further constrains the $B_\lambda$.

%======================================================================
\section{The decomposition}
%======================================================================

Everything appearing is \emph{ungerade}, and this needs no computation:
$\rho(-I)=(-I)\otimes(-I)\otimes(-I)=-\mathbb{1}_{27}$, so inversion acts as
$-1$ on all of $V$ and every gerade irreducible is excluded. Computed
multiplicities confirm it --- $A_{1g}$, $T_{1g}$, $T_{2g}$ and the
symmetric-traceless $E_g\oplus T_{2g}$ all return $0$.

The characters were built explicitly rather than copied from a table, each
verified irreducible by $\langle\chi,\chi\rangle=1$: $A_{1u}=\det R$;
$A_{2u}=\operatorname{sgn}(P)\det R$ with $P$ the permutation part;
$T_{1u}=\operatorname{tr}R$ (the defining representation);
$T_{2u}=T_{2g}\!\cdot\!\det R$ with
$T_{2g}=\tfrac12\bigl((\operatorname{tr}R)^2-\operatorname{tr}R^2\bigr)\operatorname{sgn}(P)$;
and $E_u=E_g\!\cdot\!\det R$ with $E_g$ the symmetric-traceless part less
$T_{2g}$. Then
\[
  \boxed{\;T_{1u}\otimes T_{1u}\otimes T_{1u}
   \;=\; A_{1u}\;\oplus\;A_{2u}\;\oplus\;2\,E_u\;\oplus\;4\,T_{1u}\;\oplus\;3\,T_{2u}\;}
\]

\begin{table}[h]
\centering
\begin{tabular}{lrrrr}
\toprule
irreducible & $d_\lambda$ & $m_\lambda$ & block $B_\lambda$ & copies \\
\midrule
$A_{1u}$ & 1 & 1 & $1\times1$ & 1 \\
$A_{2u}$ & 1 & 1 & $1\times1$ & 1 \\
$E_u$    & 2 & 2 & $2\times2$ & 2 \\
$T_{2u}$ & 3 & 3 & $3\times3$ & 3 \\
$T_{1u}$ & 3 & 4 & $4\times4$ & 3 \\
\midrule
 & \multicolumn{2}{r}{$\sum m_\lambda d_\lambda=27$} &
   \multicolumn{2}{r}{$\sum m_\lambda^2=31$} \\
\bottomrule
\end{tabular}
\end{table}

Both totals are independent checks and both are met: $27$ is $\dim V$, and $31$
was obtained twice by separate routes --- the character inner product
$\frac{1}{48}\sum_R(\operatorname{tr}R)^6=1488/48$, and an exact linear solve for
the commutant. The count of five irreducibles matches the dimension of the
centre of the commutant, computed independently.

\paragraph{Consistency with the eigenvalue probe.} The degeneracies
$\{1,2,3^{\times8}\}$ follow: one eigenvalue of degeneracy $1$ from each of
$A_{1u}$ and $A_{2u}$, two of degeneracy $2$ from $E_u$, and $4+3=7$ of
degeneracy $3$ from $T_{1u}$ and $T_{2u}$ --- eleven distinct eigenvalues in
all. The probe reported ten because one degeneracy-$1$ and one degeneracy-$2$
eigenvalue fell within its clustering tolerance and were counted as a single
degeneracy-$3$ value. The exact characteristic polynomial shows the same near
coincidence. This is why the decomposition is settled by characters and by the
commutant, not by a degeneracy count.

\paragraph{The practical result.}
\[
  M^{-1} \;=\; \bigoplus_\lambda \mathrm{id}_{\Gamma_\lambda}\otimes B_\lambda^{-1},
  \qquad
  \text{blocks } 1\times1,\;1\times1,\;2\times2,\;3\times3,\;4\times4 .
\]
The largest object to invert symbolically is a $4\times4$. That is closed form at
general $\lambda,\mu$ and general frequency, and it is the whole content of the
$27\times27$ inverse.

%======================================================================
\section{Why this is the right route}
%======================================================================

The practical claim is not that a symbolic inverse is unobtainable otherwise ---
one exists, at roughly a megabyte of expression. It is that such an object
answers no question. The block form answers several:

\begin{itemize}
\item each $B_\lambda^{-1}$ is small enough to read, so the frequency and
      modulus dependence of each symmetry channel is visible;
\item a channel that misbehaves (the shear channel, in this project's history)
      is isolated in one block rather than smeared over $729$ entries;
\item the static limit can be taken blockwise and checked against the
      independently derived static system;
\item the construction is valid at general $\lambda,\mu$, so it is not tied to
      the Poisson-solid values at which the existing inverse was demonstrated.
\end{itemize}

\end{document}
